
\subsection{Admissibility-Induced Topological Rigidity: A Minimal Model}
  \label{subsec:admissibility_rigidity_minimal_model}

  We present a minimal proof of principle.
  We show that a purely spectral admissibility constraint on pairs of projectors can induce a non-trivial fundamental
  group for the relative admissible space.
  This supports the claim that topological rigidity may originate from projective admissibility, not from a postulated
  ambient geometry.

  \paragraph{Spectral invariant for a pair of projectors.}
    Let $\Pi_a$ and $\Pi_b$ be rank-one projectors on $\mathbb{C}^2$.
    Consider the positive operator
    \begin{equation}
      M_{ab} = \Pi_a \Pi_b \Pi_a.
    \end{equation}
    For rank-one projectors, $M_{ab}$ has a single non-zero eigenvalue
    \begin{equation}
      \lambda_{ab} = \mathrm{Tr}(\Pi_a \Pi_b) = |\langle a|b\rangle|^2,
      \qquad
      0 \le \lambda_{ab} \le 1.
    \end{equation}
    Geometrically, writing $\Pi(\mathbf{n}) = \tfrac12(I+\mathbf{n}\cdot\boldsymbol{\sigma})$ with $\mathbf{n}\in S^2$,
    one has $\lambda_{ab}=\cos^2(\alpha/2)$, where $\alpha$ is the angle between $\mathbf{n}_a$ and $\mathbf{n}_b$.

  \paragraph{Commutativity locus and its two physical meanings.}
    The commutator norm is controlled by the same angle.
    Up to a fixed constant depending on the chosen operator norm,
    \begin{equation}
      \|[\Pi_a,\Pi_b]\| \propto \sin\alpha
      = 2\sqrt{\lambda_{ab}(1-\lambda_{ab})}.
    \end{equation}
    The commutativity locus corresponds to $\sin\alpha=0$, i.e. $\lambda_{ab}\to 1$ or $\lambda_{ab}\to 0$.
    These two limits have distinct interpretations in the cosmochrony framework.
    The limit $\lambda_{ab}\to 1$ expresses projective indistinguishability, where two descriptions collapse to the same
    effective class.
    The limit $\lambda_{ab}\to 0$ expresses complete orthogonality, corresponding to disjoint sectors with vanishing
    overlap.
    Both limits extinguish non-commutativity, but their admissibility justifications differ.

  \paragraph{Admissibility as a spectral gap on the overlap.}
    Introduce a minimal spectral separation condition on the pair,
    \begin{equation}
      \min(\lambda_{ab},1-\lambda_{ab}) \ge \delta,
      \qquad
      0 < \delta < \tfrac12.
    \end{equation}
    This excludes both near-indistinguishability and near-orthogonality,
    while enforcing a uniform lower bound on non-commutativity within the admissible region.
    Equivalently,
    \begin{equation}
      \sin\alpha \ge \varepsilon(\delta),
      \qquad
      \varepsilon(\delta)=2\sqrt{\delta(1-\delta)}.
    \end{equation}
    Here $\delta$ plays the role of a proxy for a projective gap constraint,
    and $\varepsilon$ is not ad hoc but the geometric translation of a bound on the spectral invariant $\lambda_{ab}$.

  \paragraph{Emergent non-trivial $\pi_1$ from admissibility.}
    Fix $\Pi_a$.
    The admissible set of $\Pi_b$ is then $S^2$ minus two excluded caps around $\mathbf{n}_a$ and $-\mathbf{n}_a$.
    Topologically,
    \begin{equation}
      S^2 \setminus (2\ \text{caps}) \simeq S^1 \times (0,1),
      \qquad
      \pi_1 \cong \mathbb{Z}.
    \end{equation}
    Therefore, the spectral admissibility constraint alone produces a non-trivial fundamental group for the relative
    admissible space.
    The excluded caps coincide with the locus where commutativity turns on,
    so the inhomogeneous distribution of non-commutativity is automatically realized.
    A loop that winds around the cylinder cannot be contracted without entering the spectrally inadmissible region.
    This realizes the inadmissible-border mechanism for topological rigidity.

  \paragraph{Limitations and scope.}
    This result concerns a relative configuration space (with $\Pi_a$ fixed) rather than the configuration space of two
    indiscernible excitations.
    For anyonic statistics, the relevant object is the unordered configuration space
    $\mathrm{Conf}_2(\mathcal{A})/S_2$,
    and the quotient may modify $\pi_1$.
    Moreover, the present construction is restricted to $n=2$ and does not establish Artin braid relations,
    which require a non-trivial coupling structure in the admissible space of triplets and higher $n$.
    Accordingly, this construction should be read as a proof of principle that projective admissibility can generate
    homotopic rigidity, not as a derivation of anyonic braiding.

  \paragraph{Towards $n\ge 3$ and planar versus spherical braids.}
    For $n\ge 3$, configuration spaces of projectors naturally lead to spherical braid groups when working globally on
    $\mathbb{CP}^1\simeq S^2$.
    The fundamental group of $\mathrm{Conf}_n(S^2)$ is the spherical braid group,
    which differs from the planar braid group by an additional global relation.
    This is not an obstacle for effective anyonic physics, which is planar.
    A clean route is to restrict to an admissible open set contained in a single contractible chart of $S^2$,
    away from the excluded caps.
    This requires $\delta$ sufficiently small so that the excluded caps leave a contractible region large enough to
    accommodate $n$ well-separated excitations.
    Such a condition defines an explicit regime of admissibility.
    Establishing how this regime emerges from intrinsic spectral parameters
    and how the associated configuration space supports planar braid relations
    is the next step of the program.
